\section{Заключение}

В результате выполнения учебно-исследовательской работы было выполнено следующее:

\begin{itemize}
    \item Повторён курс математической статистики в части анализа выборки, автокорреляционного анализа и построения связанных графиков.
    \item Дополнительно изучены некоторые законы распределения при помощи книг, указанных в списке литературы.
    \item Написан комплекс скриптов на языке Python, который позволяет генерировать выборку в соответствии с заданным законом распределения, рассчитывать статистические величины и создавать связанные таблицы с графиками.
\end{itemize}

После выполнения описанных действий были сделаны следующие выводы:
\begin{itemize}
    \item Исследованную выборку можно считать \textit{случайной}: она не является монотонной, не наблюдается свойство периодичности, а автокорреляционный анализ не выявил статистическую периодичность.
    \item Аппроксимация при помощи закона распределения Эрланга оказалась \textit{точной}: относительная погрешность статистических моментов не превышает $10\%$, а аналитическая плотность закона повторяет форму гистограммы для заданной выборки.  
\end{itemize}

Открытый исходный код отчёта можно изучить по ссылке ниже:
\begin{itemize}
    \item \url{https://github.com/MrDvD/itmo-labs-3/tree/main/modelling/1}
\end{itemize}

